% Transcription — fonds Alexandre Grothendieck, shelfmark 129, batch 2 (pages 21-40).
% Established from the facsimile archives/batches/129/batch-02.pdf, prepared
% from a copy of 129.pdf fetched through the project's relay
% (https://grothendieck-relay.onrender.com/source/129.pdf; PDF page k =
% archivists' page 20+k, no cover sheet), per the skill transcribe-grothendieck.
% The folder is 114 pages in six batches, transcribed in parallel by six
% passes; this is the second (pages 21-40).
%
% Pass: Opus 5.5 (claude-opus-5-5), 2026-09-26 — first pass, unchecked against the pages by a human.
%
% Dating: \dating{[1970]} is the folder's own date in src/content/catalogue.ts,
% copied verbatim; since the folder has a date of its own, the group rule
% (« Descentes », 126-133) does not apply. Nothing on these leaves dates them.
%
% Pages skipped: none.
%
% Pages identified, not transcribed (other people's typescript): all twenty.
% By the folder's rule the typed text is never re-set; each page gets one
% short French \note identifying it (subject, section numbers), followed by
% every handwritten intervention on it, given in full.
%
% Handwriting. Pages 21-28 carry none (margins and faint marks checked at
% 200-300 dpi; what shows is show-through from the facing leaf). Pages 29-40
% (chapter II) carry: a folio « II.1 » ... « II.12 » in ink, top right; and on
% pages 29-37 and 39 proof corrections — a cross (×) in the left margin
% against the line, often with the corrected symbol repeated in the margin
% ($\mu$, $\mathcal{J}$, $\mu_{n|S}$, $\mathbb{F}_p$, ↦, $\alpha_{p|S}$, $\pi$,
% $\ell$), and in the line the struck word or symbol with its replacement
% (« bigèbre » for « bi-algèbre », « directement » for « diversement »,
% « f.p.q.c.) », « = 0 », a $\pi$ for a typed p, script $\ell$ inked into
% gaps). Whether this hand is Grothendieck's cannot be told from the leaves:
% the corrections are of the kind an author or a reader makes on a
% typescript, with nothing personal in them, and the notes say so rather
% than attribute them. They are therefore given inside each page's \note,
% not as \marginal{}, which would assert they are his. One pencilled mark on
% page 29, rubbed out, is \ill{}.
%
% His pagination: none of his. The chapter II folios « II.1 »-« II.12 » are
% handwritten, hand undetermined; chapter I (pages 21-27) shows no folio.
%
% What the batch is about. (1) Pages 21-27: the end of the second typing of
% « Chapitre I. Préliminaires sur Witt » begun at page 15 (batch 1) — end
% of §3 (the ind-schemes, remark on perfect rings), §4 « L'anneau W » (W(k)
% for k perfect, the isomorphisms W/p^nW ≅ W_n, the dualising module K/W,
% Prop. 4.2 and its proof), §5.1-5.4 (absolute and relative Frobenius, the
% Verschiebung of a flat commutative group scheme after SGA 3 VII_A, their
% n-th iterates). (2) Page 28: title page of « Chapitre II. Groupes finis
% localement libres et théorie classique de Dieudonné ». (3) Pages 29-40:
% chapter II, §1 (finite locally free group schemes, augmentation ideal,
% rank, examples G_S, μ_n, α_p, Cartier duality and its exchange of
% Frobenius and Verschiebung, quotients, ℓ-primary decomposition), §2
% (infinitesimal, étale, unipotent, multiplicative, bi-infinitesimal groups;
% the Frobenius/Verschiebung criteria 2.5.1-2.5.5), §3 (the canonical
% decomposition G = G^mult × G^bi × G^et over a perfect field), and the start
% of §4 « Théorie de Dieudonné sur un corps parfait » (4.1, the Dieudonné
% ring), breaking off where the batch ends. No author's name on these pages;
% batch 1 records « DELALE » pencilled on the chapter I title page.
%
% Hardest pages: 29 (the insertion bracket and overwritten symbol in the
% bigèbre formula, a rubbed-out pencil mark) and 39 (an inked comma and
% stroke). \ill{}: 1. \uncertain{}: 1.

\documentclass[11pt,a4paper]{article}
% Le préambule commun à toutes les transcriptions du fonds Grothendieck.
%
% Il tient en une idée : **l'incertitude se compose, elle ne se résout pas.**
% Une transcription qui lisse un mot douteux a détruit la seule chose pour
% laquelle on la faisait — un lecteur doit pouvoir savoir, à chaque endroit,
% ce qui était sur le feuillet et ce qui a été ajouté. Les six macros qui
% suivent sont donc l'appareil critique, et non de la décoration.
%
% Les mêmes commandes sont reconnues par `scripts/render.mjs`, qui produit la
% vue de lecture du site. Une macro ajoutée ici doit l'être aussi là-bas,
% faute de quoi le rendu échouera bruyamment — ce qui est voulu : un rendu
% silencieusement faux est pire qu'un rendu absent.

\ProvidesPackage{grothendieck}[2026/08/08 Transcriptions du fonds Grothendieck]

\usepackage{amsmath,amssymb,amsthm}
\usepackage{mathrsfs}
\usepackage{stmaryrd}
% Les diagrammes commutatifs sont partout dans ce fonds : c'est la langue
% même de la géométrie algébrique de Grothendieck, et une transcription qui
% les rendrait en prose ne transcrirait rien.
\usepackage{tikz-cd}
% Français par défaut — c'est la langue des feuillets — mais l'anglais est
% chargé aussi : la traduction et le résumé sont des documents anglais, et la
% typographie française (espaces avant les deux-points) y serait une faute.
% Ces documents basculent par \selectlanguage{english} en tête de corps.
\usepackage[english,french]{babel}
\usepackage{fontspec}
\usepackage{geometry}
\geometry{a4paper,margin=25mm,marginparwidth=28mm}
\usepackage{marginnote}
\usepackage{xcolor}
% ulem, not soul. `soul`'s \st and \ul are built on a character-by-character
% scan that XeLaTeX's font machinery breaks: the document fails with an
% undefined control sequence rather than degrading. ulem does the same two jobs
% and is Unicode-engine safe. `normalem` stops it hijacking \emph, which the
% transcriptions use for Grothendieck's own emphasis.
\usepackage[normalem]{ulem}
\usepackage{hyperref}
\hypersetup{colorlinks=true,linkcolor=black,urlcolor=black,pdfborder={0 0 0}}

% --- Métadonnées du lot -----------------------------------------------------
% Reprises telles quelles par le script de rendu, qui les affiche en tête de
% la vue de lecture. Elles fixent ce que le fichier prétend couvrir : sans
% elles, une transcription flotte sans point d'attache dans un fonds de seize
% mille feuillets.

\newcommand{\folder}[1]{\gdef\grothFolder{#1}}
\newcommand{\batch}[1]{\gdef\grothBatch{#1}}
\newcommand{\leaves}[2]{\gdef\grothFirst{#1}\gdef\grothLast{#2}}
\newcommand{\foldertitle}[1]{\gdef\grothFoldertitle{#1}}
\gdef\grothFolder{?}\gdef\grothBatch{?}\gdef\grothFirst{?}\gdef\grothLast{?}\gdef\grothFoldertitle{}

% --- Le résumé --------------------------------------------------------------
%
% Il ouvre la lecture modernisée, sous le titre « Résumé », et il est composé
% en petit. Ce n'est pas un ornement : c'est le seul passage du document qui
% ne soit pas des mathématiques, et un lecteur doit pouvoir le reconnaître
% avant de l'avoir lu — puis le sauter s'il connaît déjà le sujet, ou s'y
% arrêter s'il ne le connaît pas. Le filet qui le suit marque la reprise du
% corps.

\newenvironment{resume}
  {\small\setlength{\parskip}{0.45em}}
  {\par\medskip\hrule\medskip}

% --- Mots-clés ---------------------------------------------------------------
%
% En anglais, en clôture du résumé de la lecture modernisée : le vocabulaire
% moderne sous lequel chercher ce que les feuillets construisent. Le manifeste
% du site (`npm run manifest`) les extrait pour étiqueter la cote entière —
% cette ligne est la source unique des tags, il n'y a pas de fichier à côté.
\newcommand{\keywords}[1]{\par\smallskip\noindent
  {\footnotesize\textsc{Keywords} --- \textit{#1}}\par}

% --- L'appareil critique ----------------------------------------------------

% Le numéro de feuillet, en marge. C'est l'ancre qui permet de régler un
% désaccord sur une formule en regardant *un* feuillet plutôt qu'une chemise
% de six cents. Le site s'en sert aussi pour tourner les pages du fac-similé
% quand on fait défiler la transcription.
\newcommand{\leaf}[1]{%
  \marginnote{\footnotesize\color{gray!70}\textsf{#1}}%
  \label{leaf:#1}\ignorespaces}

% Illisible : on ne devine pas. Un mot inventé qui se lit comme les autres est
% le pire résultat possible.
\newcommand{\ill}{\textcolor{red!60!black}{[\,\ldots\,]}}

% Lecture douteuse : le mot est donné, et signalé comme incertain.
\newcommand{\uncertain}[1]{\uline{#1}}

% Ajout de l'éditeur — une lettre manquante, un mot sous-entendu.
\newcommand{\add}[1]{\textcolor{blue!50!black}{[#1]}}

% Biffé par Grothendieck lui-même. Ce qu'il a rayé fait partie du document :
% une rature dit souvent où la pensée a changé de direction.
\newcommand{\struck}[1]{\sout{#1}}

% Note du transcripteur, sur l'état matériel du feuillet.
\newcommand{\note}[1]{\par\smallskip{\small\color{gray!60!black}\textit{[#1]}}\par\smallskip}

% Note marginale de Grothendieck, distincte du corps du texte.
\newcommand{\marginal}[1]{\par\smallskip\noindent\hspace{1em}%
  {\small\color{gray!60!black}#1}\par\smallskip}

% --- Titre ------------------------------------------------------------------

\AtBeginDocument{%
  \begin{center}
    {\large\bfseries Fonds Alexandre Grothendieck --- cote n\textsuperscript{o}~\grothFolder}\\[2pt]
    {\itshape\grothFoldertitle}\\[4pt]
    {\small feuillets \grothFirst--\grothLast\ (lot \grothBatch)}\\[2pt]
    {\footnotesize\color{gray!60!black}Transcription établie d'après le fac-similé numérisé
     par l'Université de Montpellier.\\
     Les lectures incertaines sont soulignées, les illisibles notées [\,\ldots\,],
     les ajouts entre crochets.}
  \end{center}
  \vspace{1em}}

\endinput


\folder{129}
\batch{2}
\pages{21}{40}
\dating{[1970]}
\watermark{Édition de démonstration}
\foldertitle{Notes séminaire Montréal (Delale, Labute, Hakim, Messing) : copies de tapuscrits annotés (s.d.), notes manuscrites (s.d.)}

\begin{document}

\page{21}
\note{seconde frappe du tapuscrit « Chapitre I. Préliminaires sur Witt »
(attribué à Delale, cf.\ page 1), suite de la page 20 : fin du §3
(multiplication par $p$ comme morphisme de transition, Remarque sur un
anneau parfait) et §4 « L'anneau $W$ », n°~4.1 ($W(k)$ pour $k$ parfait,
renvoi à Serre, \emph{Corps locaux}). Texte d'autrui, non transcrit ;
aucune annotation manuscrite}

\page{22}
\note{même tapuscrit, suite du n°~4.1 (isomorphismes
$\varepsilon_n : W/p^nW \simeq W_n$, diagramme commutatif, module
dualisant $K/W$) et énoncé de la Proposition~4.2. Aucune annotation
manuscrite}

\page{23}
\note{même tapuscrit, démonstration de la Proposition~4.2. Aucune
annotation manuscrite}

\page{24}
\note{même tapuscrit, fin de la démonstration de 4.2, Remarque, et début du
§5 « Les morphismes de Frobenius et de Verschiebung d'un schéma en
groupes », n°~5.1 (Frobenius absolu). Aucune annotation manuscrite}

\page{25}
\note{même tapuscrit, fin du n°~5.1 et n°~5.2 (Frobenius relatif
$X \to X^{(p/S)}$). Aucune annotation manuscrite}

\page{26}
\note{même tapuscrit, fin du n°~5.2 et n°~5.3 (Verschiebung d'un schéma en
groupes commutatifs plat, renvoi à SGA~3, VII$_{\mathrm{A}}$, 4.2--4.3).
Aucune annotation manuscrite}

\page{27}
\note{même tapuscrit, n°~5.4 (itérés du Frobenius relatif et de la
Verschiebung) ; fin du chapitre~I. Aucune annotation manuscrite}

\page{28}
\note{page de titre dactylographiée : « Chapitre II. Groupes finis
localement libres et théorie classique de Dieudonné ». Aucun nom d'auteur ;
aucune intervention manuscrite}

\page{29}
\note{tapuscrit, chapitre~II, §1 « Rappels généraux sur les groupes finis
localement libres », n°~1.1 (algèbre affine $A$, morphismes
$\Delta$, $I$, $\varepsilon$, bigèbre, idéal d'augmentation) et début du
n°~1.2 (rang). Texte d'autrui, non transcrit. Interventions manuscrites,
d'une main qu'on ne peut pas identifier d'après le feuillet (rien ne
permet de dire si c'est celle de Grothendieck) : en haut à droite, le folio
«~II.1~» ; une accolade à l'encre part d'après « muni de sa loi
d'algèbre » et embrasse la formule détachée, dont le symbole
dactylographié devant le deux-points est surchargé en $\mu$, ce qui donne
$\mu : A \otimes A \to A$ ; en marge, une croix et une petite coche. Dans
« Le noyau J de $\varepsilon$ », le J dactylographié est biffé de deux
traits, et la marge porte un $\mathcal{J}$ cursif. Plus bas dans la marge,
une inscription au crayon, gommée : \ill{}}

\page{30}
\note{même tapuscrit, fin du n°~1.2 (groupe de rang $n$) et n°~1.3
« Exemples de groupes finis localement libres », 1.3.1--1.3.3 (groupe
constant $G_S$, groupe multiplicatif et $\mu_{n|S}$, groupe additif).
Interventions manuscrites, main non identifiée : folio «~II.2~» ; devant
« Par définition, $\mu_{n|S}$ », le symbole dactylographié est barré, et
la marge porte une croix et $\mu_{n|S}$, le $\mu$ à hampe doublée comme
une lettre ajourée ; devant la ligne « Si $S$ est un schéma sur le corps
premier $\mathbb{F}_p$ », une croix et $\mathbb{F}_p$ en marge, sans
correction visible dans la ligne}

\page{31}
\note{même tapuscrit, fin du n°~1.3.3 ($F$, $\alpha_{p|S}$) et n°~1.4
« Dualité de Cartier » ($G^* = \underline{\mathrm{Hom}}(G,
\mathbb{G}_{m|S})$, bidualité, renvoi à SGA~3, VII$_{\mathrm{A}}$,
3.3.1). Interventions manuscrites, main non identifiée : folio «~II.3~» ;
en marge, une croix et le signe $\mapsto$, devant « foncteur contravariant
$G \mapsto G^*$ » ; en marge une croix, et dans la ligne « bi-algèbre »
biffé, remplacé au-dessus par « bigèbre » ; en marge une croix, et dans
« $G \to (G^*)^*$ » l'astérisque de l'exposant ajouté à l'encre}

\page{32}
\note{même tapuscrit, suite du n°~1.4 (échange du Frobenius relatif et de
la Verschiebung, renvoi à SGA~3, VII$_{\mathrm{A}}$, 4.3.3), n°~1.5
« Exemples » ($(\mathbb{Z}/n\mathbb{Z})^*_S \simeq \mu_{n|S}$,
$(\alpha_{p|S})^* \simeq \alpha_{p|S}$) et début du n°~1.6 « Quotients de
groupes finis localement libres ». Interventions manuscrites, main non
identifiée : folio «~II.4~» ; en marge une croix, et dans « qui résulte
aussi diversement des définitions » le mot « diversement » biffé, remplacé
au-dessus par « directement » ; en marge une croix et $\alpha_{p|S}$,
devant $(\alpha_{p|S})^* \simeq \alpha_{p|S}$ ; en marge $\pi$, et dans la
suite exacte $0 \to G' \to G \to G'' \to 0$ le p au-dessus de la seconde
flèche surchargé en $\pi$}

\page{33}
\note{même tapuscrit, suite du n°~1.6 (représentabilité de $G''$, renvoi à
SGA~3, V, 4.1 ; multiplicativité du rang) et n°~1.7 (composante
$\ell$-primaire, renvoi à SGA~3, VIII, 7.3). Interventions manuscrites,
main non identifiée : folio «~II.5~» ; en marge $\pi$, et dans « et que p
est fidèlement plat » le p biffé d'une croix, $\pi$ écrit au-dessus ;
deux fois, en marge, une croix et $\ell$, devant « de $\ell$-torsion » et
« un $\ell$-groupe », où le $\ell$ est écrit à la main dans le blanc laissé
par la frappe}

\page{34}
\note{même tapuscrit, fin du n°~1.7 (groupes $p$-primaires, dual d'un
$p$-groupe) et §2 « Types particuliers de groupes », n°~2.1, 2.1.1
(infinitésimal) et 2.1.2 (étale). Interventions manuscrites, main non
identifiée : folio «~II.6~» ; en marge une croix, et en bout de la ligne
« (pour la topologie étale, ou f.p.p.f., ou » l'ajout « f.p.q.c.) »}

\page{35}
\note{même tapuscrit, 2.1.3 (unipotent, renvoi à SGA~3, XVII), 2.1.4 (type
multiplicatif), 2.1.5 (bi-infinitésimal), n°~2.2 (en caractéristique $p$,
un $p$-groupe de type multiplicatif est infinitésimal) et début du n°~2.3.
Interventions manuscrites, main non identifiée : folio «~II.7~» ; une
petite croix en marge devant 2.1.5, sans correction visible dans la
ligne}

\page{36}
\note{même tapuscrit, fin du n°~2.3 et n°~2.4 « Exemples de $p$-groupes en
caractéristique $p$ » ($(\mathbb{Z}/p\mathbb{Z})_S$, $\mu_{p|S}$,
$\alpha_{p|S}$ ; suite de composition, renvoi à SGA~3, XVII, 1.7 et
4.2.1). Interventions manuscrites, main non identifiée : folio «~II.8~» ;
en marge une croix, et en bout de la suite de composition
$G \supset G_0 \supset \dots \supset G_n$ l'ajout « $= 0$ »}

\page{37}
\note{même tapuscrit, n°~2.5 (critères 2.5.1--2.5.5 par le Frobenius
relatif et la Verschiebung) et démonstration de 2.5.1 (idéaux $J_n$ et
puissances $J^{p^n}$). Interventions manuscrites, main non identifiée :
folio «~II.9~» ; une croix en marge devant la ligne « lons caractériser le
"noyau" de $F^n_{G/S}$ pour tout S-schéma $G$ muni d'une », sans
correction visible dans la ligne}

\page{38}
\note{même tapuscrit, fin de la démonstration de 2.5.1, démonstration de
2.5.4 (renvoi à SGA~5, XIV (Houzel)), §3 « Décomposition d'un $p$-groupe
fini localement libre sur une base de caractéristique $p$ réduite à un
point », n°~3.1 (suite exacte $0 \to G^0 \to G \to G^{\mathrm{et}} \to 0$,
sous-groupe $G^{\mathrm{mult}}$). Seule intervention manuscrite : le folio
«~II.10~»}

\page{39}
\note{même tapuscrit, suite du n°~3.1 et n°~3.2 (décomposition canonique
$G = G^{\mathrm{mult}} \times G^{\mathrm{bi}} \times G^{\mathrm{et}}$ sur un
corps parfait, renvoi à SGA~3, XVII, 1.6). Interventions manuscrites, main
non identifiée : folio «~II.11~» ; en marge une croix, et dans « est
bi-infinitésimal car c'est » une virgule ajoutée à l'encre après
« bi-infinitésimal », avec un trait appuyé devant « car »
(\uncertain{suppression d'un signe dactylographié})}

\page{40}
\note{même tapuscrit, liste des équivalences (infinitésimal, étale,
unipotent, bi-infinitésimal, multiplicatif), n°~3.3 ($p$-groupes étales
et de type multiplicatif sur un corps parfait, $\pi_1(k)$, dualité de
Pontryagin) et début du §4 « Théorie de Dieudonné sur un corps parfait »,
n°~4.1 (catégorie $p$-Gr-fin$(k)$, anneau de Dieudonné $W[F,V]$) ; la page
s'arrête en cours de phrase. Seule intervention manuscrite : le folio
«~II.12~»}

\end{document}
